\documentclass[11pt]{article}
\usepackage[margin=2.0cm]{geometry}
\usepackage{booktabs}
\usepackage{amsmath}
\usepackage{amssymb}
\usepackage{enumitem}
\usepackage{array}
\usepackage{listings}
\usepackage{xcolor}
\usepackage[T1]{fontenc}
\usepackage{lmodern}

\setlength{\parindent}{0pt}
\setlength{\parskip}{3.6pt}
\newcommand{\code}[1]{\texttt{#1}}
\newcommand{\hr}{\vspace{2pt}\hrule\vspace{6pt}}
\renewcommand{\arraystretch}{1.12}

\definecolor{kw}{RGB}{0,90,160}
\definecolor{cmt}{RGB}{110,110,110}
\lstdefinestyle{vhdl}{
  language=VHDL, basicstyle=\ttfamily\footnotesize,
  keywordstyle=\color{kw}\bfseries, commentstyle=\color{cmt}\itshape,
  showstringspaces=false, columns=fullflexible, breaklines=true,
  frame=leftline, framesep=6pt, xleftmargin=8pt, aboveskip=6pt, belowskip=6pt
}
\lstset{style=vhdl}

\begin{document}

\begin{center}
{\Large\bfseries EE\,362 --- Computer Organization Laboratory}\\[2pt]
{\large\bfseries Lab 1 --- Solution}\\[2pt]
{\itshape The Transparent Latch}\\[4pt]
Fall 2026 \quad$\cdot$\quad Lab TA: Gandhi Torres
\end{center}
\hr

Worked solution for \code{TransLatch} and \code{TransLatch8}. Read it after you
have made your own attempt run --- the two bugs in Section 5 are the ones this lab
was built to give you, and reading about them is a poor substitute for hitting them.

\vspace{4pt}\hr

\section*{1. Specification}

A one-bit storage element:

\begin{itemize}[leftmargin=1.6em,itemsep=1pt]
\item \code{Clr}\;=\;1 forces \code{Q} to 0 \textbf{immediately}, whatever else is
happening. Asynchronous: it waits for nothing.
\item \code{Clr}\;=\;0, \code{En}\;=\;1: the latch is \textbf{transparent}. \code{Q}
follows \code{D} continuously, not on an edge.
\item \code{Clr}\;=\;0, \code{En}\;=\;0: the latch \textbf{holds}. \code{Q} keeps
the value it had when \code{En} went low, and \code{D} is ignored.
\end{itemize}

The third bullet is the entire lab. It is the first circuit in this course whose
output depends on something other than its present inputs.

\section*{2. Truth table --- Type T (transition)}

A combinational table maps inputs to outputs. A \textbf{transition table} maps
inputs \emph{and the present state} to the \emph{next} state. The column is written
$Q^{+}$, read ``Q next''.

\begin{center}
\begin{tabular}{c|ccc|c|l}
\toprule
\# & \code{Clr} & \code{En} & \code{D} & $Q^{+}$ & note \\
\midrule
0 & 1 & -- & -- & \textbf{0} & asynchronous clear --- dominates everything \\
1 & 0 & 0 & -- & $\mathbf{Q}$ & \textbf{hold} --- the memory row \\
2 & 0 & 1 & 0 & \textbf{0} & transparent \\
3 & 0 & 1 & 1 & \textbf{1} & transparent \\
\bottomrule
\end{tabular}
\end{center}

Three things to notice, because each one is a rule you will use for the rest of the
course.

\textbf{Row 1 has $Q$ in the output column.} Not 0, not 1 --- $Q$, the value it
already had. \textbf{Any output column containing $Q$ means the design has memory.}
That is the signature you look for, and from Lab 2 onwards you will be reading it
rather than writing it.

\textbf{Row 0 comes first, and its don't-cares are a priority claim.} Writing
\code{--} under \code{En} and \code{D} asserts that when \code{Clr} is high, nothing
else can change the outcome. Row 0 is listed above rows 2 and 3 because in a table
with overlapping rows, \textbf{the first matching row wins}. Your VHDL must encode
that same priority --- and it does so by testing \code{Clr} first.

\textbf{Rows 2 and 3 could be compressed to one row reading $Q^{+} = D$.} They are
written out because at four rows it costs nothing, and because seeing $D$ appear in
the output column next to $Q$ makes the contrast between ``pass'' and ``hold''
visible at a glance.

\section*{3. Boolean derivation}

\textbf{(1) Canonical --- read straight off the table.} Take the rows where
$Q^{+}$ can be 1, and write the product of conditions that reaches each:

\[
Q^{+} = \underbrace{\overline{\code{Clr}} \cdot \code{En} \cdot D}_{\text{row 3: pass a 1}}
\;+\;
\underbrace{\overline{\code{Clr}} \cdot \overline{\code{En}} \cdot Q}_{\text{row 1: hold a 1}}
\]

Row 0 contributes nothing (its output is 0) and row 2 contributes nothing (its
output is 0). Row 1 is the unusual one: its term contains $Q$, the output itself.
\textbf{That feedback is the memory.} An equation whose right-hand side mentions its
own left-hand side is not a mistake here --- it is the definition of state.

\textbf{(2) Simplified --- factor out $\overline{\code{Clr}}$:}
\[
Q^{+} = \boxed{\overline{\code{Clr}} \cdot \left( \code{En} \cdot D + \overline{\code{En}} \cdot Q \right)}
\qquad \textbf{[distributive law]}
\]

Now read what the parenthesis is: $\code{En}\cdot D + \overline{\code{En}}\cdot Q$ is
exactly the 2-to-1 multiplexer from Lab 0 Exercise 2, with $\code{En}$ as the select,
$D$ as one input --- and $Q$, the latch's own output, as the other.

\begin{quote}
\textbf{A transparent latch is a multiplexer with its output fed back to one of its
inputs.} That is all it is. When \code{En} selects $D$, new data flows in; when
\code{En} selects $Q$, the old value circulates and sustains itself.
\end{quote}

The $\overline{\code{Clr}}$ factored outside the parenthesis is the priority from
row 0, made structural: an AND gate on the output that can force 0 regardless of
what the multiplexer chose.

\textbf{(3) VHDL.}
\begin{lstlisting}
process (Clr, En, D)
begin
    if Clr = '1' then
        Q <= '0';        -- row 0: asynchronous clear wins
    elsif En = '1' then
        Q <= D;          -- rows 2, 3: transparent
    end if;
    -- row 1 (hold) is the deliberately missing else
end process;
\end{lstlisting}

\subsection*{3.1 The missing \code{else} --- read this twice}

There is no \code{else}. That is not an omission; it is the implementation of row 1.

The VHDL rule is:
\begin{quote}
\textbf{In a process, a signal that is not assigned on some path keeps its previous
value.}
\end{quote}

When \code{Clr} = 0 and \code{En} = 0, neither branch executes, \code{Q} is never
assigned, and so \code{Q} holds. Synthesis reads that requirement --- ``this signal
must remember its value'' --- and gives you a latch. Which is exactly what you
asked for.

Everywhere else in your career this is a bug with a name: an \textbf{inferred
latch}, the thing your synthesis tool warns you about when you forget a branch in a
combinational \code{case}. \textbf{Lab 1 is the one time all semester you want it.}
Lab 3 is the first time you do not, and by then you will recognise the warning
because you deliberately caused it here.

\section*{4. TransLatch8 --- the same equation, subscripted}

\[
Q_i^{+} = \overline{\code{Clr}} \cdot \left( \code{En} \cdot D_i + \overline{\code{En}} \cdot Q_i \right)
\qquad i = 0 \ldots 7
\]

Eight independent copies of the same circuit, sharing two control signals. \code{Clr}
and \code{En} stay \textbf{one bit wide} --- they are broadcast to all eight bits.

\begin{quote}
\textbf{Control signals are narrow; datapaths are wide.} This is the structural
pattern of the entire course. In Lab 5 the ALU is 32 bits wide with a 4-bit function
code. In Lab 12 the CPU is 32 bits wide with a 16-field control record. It starts
here, at 8 and 1.
\end{quote}

The code changes in exactly two places: the port types, and the cleared value.

\begin{lstlisting}
process (Clr, En, D)
begin
    if Clr = '1' then
        Q <= (others => '0');   -- row 0
    elsif En = '1' then
        Q <= D;                 -- rows 2, 3
    end if;
    -- row 1 (hold): no else
end process;
\end{lstlisting}

\code{Q <= '0';} would not compile --- \code{'0'} is a single bit and \code{Q} is a
vector. \code{(others => '0')} means ``every element gets \code{'0'}'' and is
width-independent: the identical line works at 8 bits, at 32, and at 1024. Use it
from now on, and never write \code{"00000000"}.

\section*{5. The two bugs this lab is built around}

\subsection*{5.1 The incomplete sensitivity list}

The starter file gives you \code{process (En)}. It compiles without an error. It
simulates \textbf{wrong}.

A process wakes up only when a signal in its sensitivity list changes. With only
\code{En} listed, changing \code{D} while the latch is transparent does not wake the
process, so \code{Q} does not follow \code{D} --- and changing \code{Clr} does
nothing at all until \code{En} happens to move. You get a circuit that is
transparent only at the instant \code{En} changes, which is not a latch and is not
anything.

\textbf{The rule: a process's sensitivity list must name every signal the process
reads.} Here that is \code{Clr}, \code{En} and \code{D}. Not \code{Q} --- \code{Q}
is written, not read.

\textbf{How the failure looks in simulation:} the \emph{first} check fails, with
\code{Q = 'U'}. Vector 1 raises \code{Clr} while \code{En} does not change, so the
process never wakes and nothing has ever been assigned. It looks as though no code was
written at all. Measured against the reference testbench, the 1-bit latch fails checks
1, 3, 6, 7 and 9: vector 3 because \code{D} changes while \code{En} is already high,
and the clear checks because a change of \code{Clr} alone wakes nothing.

\textbf{Which signal is missing is named by which checks fail:} clear checks failing
means \code{Clr} is missing; the ``tracks a change of D'' check failing means \code{D}
is missing. Note that on the 1-bit latch the hold checks (4 and 5) still \emph{pass} ---
a value that is stuck looks exactly like a value that is held. A test can pass for the
wrong reason. (The 8-bit latch also fails 4 and 5, because its data changed while
\code{En} stayed high and the stale value is what it holds.)

\textbf{How to see it without simulating:} Quartus reports it, once per missing
signal:
\begin{lstlisting}
Warning (10492): VHDL Process Statement warning at TransLatch.vhd(33): signal "Clr"
  is read inside the Process Statement but isn't in the Process Statement's
  sensitivity list
\end{lstlisting}
and then builds the complete latch anyway (verified: 1 inferred latch, 0 registers).
\textbf{Note what that means --- the synthesised hardware is correct and the
simulation is wrong.} They disagree.

\textbf{This depends on the \code{Clr} branch.} Remove it, as in Example 3, and the
same missing name makes Quartus build a \emph{flip-flop} clocked by \code{En}, with no
warning at all. The examples solution handout, Section 3.4, has both cases measured. That divergence,
where your simulation and your chip do different things, is the most expensive class
of bug in digital design, and you have now met it in week one on a four-row table.

\subsection*{5.2 Writing the \code{else}}

The other half of the class writes:

\begin{lstlisting}
    if Clr = '1' then  Q <= '0';
    elsif En = '1' then  Q <= D;
    else  Q <= '0';        -- <-- WRONG
    end if;
\end{lstlisting}

This is complete, warning-free, and not a latch. It is a combinational AND gate:
$Q = \overline{\code{Clr}} \cdot \code{En} \cdot D$. Row 1 of the table now reads 0
instead of $Q$.

\textbf{How it fails:} only at the two checks made while the latch is closed and
holding a 1: vector 4 (``value captured before closing'') and vector 5 (the hold
vector). Every check made with the latch open, or after a clear, still passes. A
design that passes eight of ten checks and fails exactly the ones that define the
component is a much better lesson than one that fails everywhere.

\section*{6. The testbench --- why these ten vectors}

Ten vectors, each checked on both the 1-bit and the 8-bit latch: 20 assertions.

\begin{center}
\begin{tabular}{clll}
\toprule
\# & what it drives & covers & catches \\
\midrule
1 & \code{Clr}=1 with \code{D}=1 & row 0 & \textbf{\code{Clr} missing from list} (\code{'U'}) \\
2 & transparent, \code{D}=1 & row 3 & \\
3 & change \code{D}, still open & rows 2/3 & \textbf{\code{D} missing from list} \\
4 & load, then close \code{En} & rows 3$\to$1 & \textbf{the written \code{else}} \\
5 & change \code{D} while closed & \textbf{row 1} & \textbf{the written \code{else}}, a wire \\
6 & \code{Clr}=1 while closed & row 0 over row 1 & clear gated by \code{En} \\
7 & release \code{Clr}, still closed & row 1 & a clear that ``un-clears'' \\
8 & reopen \code{En} & rows 1$\to$3 & \\
9 & \code{Clr}=1 while transparent & row 0 over row 3 & priority written backwards \\
10 & release \code{Clr}, still open & row 3 & \\
\bottomrule
\end{tabular}
\end{center}

\textbf{The coverage argument} --- this is what your report's Section 5 must contain,
in your own words:

\begin{quote}
All four table rows are exercised. Both entries into the hold state are tested
(vector 4 from transparent, vector 7 from cleared), and the hold is verified by
changing \code{D} underneath it (vector 5) rather than merely by leaving it alone ---
a latch stuck at its value would pass a passive check. The clear is asserted from
each of the other two states (vector 6 from hold, vector 9 from transparent) because
row 0's don't-cares claim it dominates both, and a claim in the table must be a
check in the testbench. The 8-bit checks use \code{A5}, \code{5A} and \code{3C},
which are bitwise complements and rotations of each other, so a swapped or
sign-extended bit ordering produces a visibly wrong value rather than an accidental
match.
\end{quote}

Note vectors 5 and 7 together. They apply \textbf{identical inputs} ---
\code{Clr}=0, \code{En}=0, \code{D}=0 --- yet expect \code{Q}=1 at vector 5 and
\code{Q}=0 at vector 7. No combinational circuit can give two different outputs for
the same inputs, so \textbf{this pair is the proof that the design has memory}, and
the coverage argument has to name it.

\textbf{Expected transcript on success:}
\begin{lstlisting}
# ** Note: PASS - TransLatch: all 20 checks passed.
\end{lstlisting}

If you see no \code{PASS} line and no \code{FAIL} lines, your simulation did not run
the process --- check that the top-level entity for simulation is \code{tb\_TransLatch}
and not \code{TransLatch}.

\section*{7. Quartus --- what to do, and what to read}

The project is supplied. \code{Lab1.qpf} already sets family Cyclone~V, device
\code{5CSEMA4U23C6}, simulator \code{Questa Altera FPGA (VHDL)}, and the testbench
\code{tb\_TransLatch}. \textbf{Open it, do not create a new one.}

\begin{enumerate}[leftmargin=2em,itemsep=2pt]
\item \code{File $\rightarrow$ Open Project} $\rightarrow$ \code{Lab1.qpf}
\item Edit \code{TransLatch.vhd}: fix the sensitivity list, implement the table.
\item \code{Processing $\rightarrow$ Start $\rightarrow$ Start Analysis \& Elaboration}
\item \code{Tools $\rightarrow$ Netlist Viewers $\rightarrow$ RTL Viewer}
\item \code{Processing $\rightarrow$ Start Compilation} \quad(Ctrl+L)
\item \code{Tools $\rightarrow$ Run Simulation Tool $\rightarrow$ RTL Simulation}
\item Repeat 2--6 for \code{TransLatch8}: \code{Project $\rightarrow$ Set as
Top-Level Entity} with \code{TransLatch8.vhd} open, then recompile.
\end{enumerate}

\subsection*{Evidence --- and what it should say}

Full descriptions of each source are in the \emph{Quartus Results Guide}. What
matters for \emph{this} design:

\vspace{4pt}
\textbf{E1 --- Flow Summary / Messages.}
\begin{center}
\begin{tabular}{ll}
\toprule
\textbf{Total registers} & \textbf{0} \\
Total pins & 4 (\code{TransLatch}) \quad / \quad 18 (\code{TransLatch8}) \\
\bottomrule
\end{tabular}
\end{center}

\textbf{Zero registers is the correct and expected result}, and it is the single most
informative number in this report. A latch is not a register --- there is no clock
anywhere in this design. If your report says ``8 registers'', you built a flip-flop,
which is Lab 2.

Pin count: 18 for the 8-bit version --- 8 in, 8 out, plus \code{En} and \code{Clr}.
The two control bits serve all eight data bits, which is that narrow/wide pattern
showing up as an actual number for the first time.

\vspace{4pt}
\textbf{E2 --- RTL Viewer.} Look for the \textbf{feedback loop}: the output of the
selection logic routed back to its own input. You are looking at Section~3's
equation drawn as a picture --- a multiplexer eating its own output, with
\code{Clr} gating the result. If you see no loop, you wrote the \code{else}.

\vspace{4pt}
\textbf{E3 --- Technology Map Viewer, and the latch report.} Open the Compilation
Report and find \textbf{Analysis \& Synthesis $\rightarrow$ User-Specified and
Inferred Latches}. For \code{TransLatch8} the reference build reports:

\begin{center}
\begin{tabular}{lll}
\toprule
Latch Name & Latch Enable Signal & Free of Timing Hazards \\
\midrule
\code{Q[0]\$latch} & \code{En} & yes \\
\multicolumn{3}{c}{\ldots\ through \ldots} \\
\code{Q[7]\$latch} & \code{En} & yes \\
\midrule
\multicolumn{3}{l}{Number of user-specified and inferred latches = \textbf{8}} \\
\bottomrule
\end{tabular}
\end{center}

Three things to take from that table, and all three belong in your report:

\begin{enumerate}[leftmargin=2em,itemsep=1pt]
\item \textbf{Eight latches for eight bits.} One storage element per bit, exactly as
the subscripted equation predicted.
\item \textbf{The enable signal is \code{En}} --- Quartus worked out on its own which
signal controls the storage, and it agrees with your table. If it names something
else, your priority order is wrong.
\item \textbf{``Free of timing hazards: yes''} means the tool checked for a glitch on
the enable path that could corrupt the stored value, and found none.
\end{enumerate}

\textbf{This report exists to warn you.} Ninety-nine times out of a hundred, latches
appear here because someone forgot a branch. This is the one lab where finding your
component listed means you got it right --- and having seen it deliberately, you will
recognise it instantly when it is accidental.

\vspace{4pt}
\textbf{E4 --- Resource utilisation.} Around 5 ALMs for the 8-bit version, under 1\%
of the device. The absolute number is uninteresting; the \emph{ratio} is not. Record
it, because you will build the same table for every lab and in Lab 12 you will be
able to say which component actually costs anything. (Spoiler for December: it is the
register file, and it is not close.)

\vspace{4pt}
\textbf{E5 --- Pin-Out file.} Skim it. Nothing is pin-constrained this semester, so
the fitter placed everything itself; the file's purpose in Lab 1 is to show you it
exists and what it looks like.

\vspace{4pt}
\textbf{E6 --- Timing Analyzer.} With no \code{.sdc} file, Quartus \textbf{invents a
clock} on the latch enable and reports no $F_{max}$:
\begin{lstlisting}
Info (332105): create_clock -period 1.000 -name En En
Info (332140): No fmax paths to report
\end{lstlisting}
There is no $F_{max}$ because there are no setup or hold paths: nothing travels from
one storage element to another. The only check left is \emph{minimum pulse width}
on \code{En} against that invented 1~ns (1~GHz) period. It may pass or fail: the
reference \code{TransLatch8} build fails it by 38~ps at the fast corner and prints
\code{Critical Warning (332148): Timing requirements not met}.

\textbf{That warning is not a design error.} It is a check against a clock nobody
specified. In your report, say that the enable was treated as a 1~ns clock by
default, that the design has no register-to-register path and so no $F_{max}$, and
that a meaningful timing number needs a real clock and a constraint. Both arrive in
Lab 2.

\vspace{4pt}
\textbf{E7 --- Questa.} The waveform, plus the transcript line proving the checks
ran. On the waveform, find and annotate one specific place: \textbf{the interval
where \code{En} is low, \code{D} changes, and \code{Q} does not.} A screenshot with
that interval marked is worth more than three paragraphs of description, and it is
the picture your Discussion section should be built around.

\section*{8. Report checklist --- \code{lab01.tex}}

\begin{itemize}[leftmargin=1.6em,itemsep=1pt]
\item Table T with rows numbered, and $Q$ (not 0) in row 1's output column
\item The three-line derivation, with \textbf{distributive law} named at step 2
\item One sentence identifying the parenthesis as a 2:1 multiplexer with feedback ---
tie it back to Lab 0 Exercise 2 explicitly
\item Both files via \code{\textbackslash lstinputlisting}, never pasted
\item Coverage argument naming \textbf{vectors 5 and 7}: same inputs, different
expected outputs, so only memory can pass both
\item Evidence: Flow Summary showing \textbf{0 registers}, the inferred-latch table,
the RTL Viewer feedback path, the annotated waveform interval
\item Two sentences on the Timing Analyzer: the invented 1~ns clock on \code{En}, and
why there is no $F_{max}$
\item Discussion: which of the two bugs you hit, which vector caught it, and how long
it took you to find. \textbf{Be honest here} --- it is graded on the quality of the
account, not on having got it right first time.
\end{itemize}

\vspace{6pt}
\hr
{\small \textbf{Forward.} Lab 2 wires two of these in series on opposite clock
phases, and the pair becomes a master--slave flip-flop --- edge-triggered storage,
and the first $F_{max}$ number of the semester. Lab 6 instantiates 1024 of them as a
register file. The transparent latch is the only storage primitive in this course;
everything after this is built from it.}

\end{document}
